\emph{Notation.} $h = L^{2}(\R, du)\otimes K$ (complex unit $I_{4}$, the field theory's unit; numbers, analytic continuations included, are read in $\mathrm{span}\{1, I_{4}\}$, with $\mathrm{Im}$ the $I_{4}$-coefficient), $h_{a} = L^{2}((a, \infty))\otimes K$, $h_{a}^{c} = L^{2}((-\infty, a))\otimes K$; $\psi$ is antilinear and $\psi^{*}$ linear in the test function, $\omega(\psi^{*}(f)\psi(g)) = \ip{g}{Tf}$, so $\omega(\psi(f)\psi^{*}(g)) = \ip{f}{(1 - T)g}$ and, $\omega$ being gauge-invariant, $\omega(B(f)B(g)) = \ip{f}{(1 - T)g} + \ip{g}{Tf}$. In the GNS representation $(\pi, \mathcal{H}, \Omega)$ the grading is implemented by a self-adjoint unitary $\Theta$ with $\Theta\Omega = \Omega$ ($\omega$ is even), $Z = (1 + I_{4}\Theta)/(1 + I_{4})$, and for an odd operator $Y$, $ZYZ^{*} = I_{4}\Theta Y$; on odd vectors $Z = -I_{4}$. $\Theta$ preserves $\mathcal{R}_{a}$ and fixes $\Omega$, so it commutes with $J_{a}$ and $\Delta_{a}$, and $J_{a}ZJ_{a} = (1 - I_{4}\Theta)/(1 - I_{4}) = Z^{*}$. Fourier convention $\hat x(p) = \int x(u)e^{I_{4}pu}\,du$, so $P$ acts as multiplication by $p$ and $\ip{f}{g} = \tfrac{1}{2\pi}\int \ip{\hat f(p)}{\hat g(p)}_{K}\,dp$; $\widehat{\rho_{0}x}(p) = \hat x(-p)$.

\emph{Step 0 (natural-cone positivity).} For $A = A^{*} \in \mathcal{R}_{a}$, Tomita's operator gives $J_{a}\Delta_{a}^{1/2}A\Omega = A\Omega$, so $\Delta_{a}^{1/2}A\Omega = J_{a}A\Omega$ and $\ip{A\Omega}{J_{a}A\Omega} = \|\Delta_{a}^{1/4}A\Omega\|^{2} \ge 0$. Only $S = J\Delta^{1/2}$ is used; no property of the modular group is.

\emph{Step 1 (what CGMA$_{\eta}$ at one cut says about the state).} Let $\tilde J_{a} = Z^{*}J_{a}$, antiunitary with $\tilde J_{a}\Omega = \Omega$. The hypothesis is $\tilde J_{a}\pi(\psi(f))\tilde J_{a}^{-1} = \eta\,\pi(\psi^{*}(\rho_{a}f))$ for $f \in h_{a}$. Applying $J_{a}(\cdot)J_{a}$ and using $J_{a}^{2} = 1$, $J_{a}ZJ_{a} = Z^{*}$ gives the same identity for $f \in h_{a}^{c}$ (with $\rho_{a}f \in h_{a}$), hence, by linearity, for all $f \in h$: $\mathrm{Ad}\,\tilde J_{a}$ restricts to the antilinear $*$-automorphism $\alpha_{a}$ of $\pi(\CAR(h))$ with $\alpha_{a}(\psi(f)) = \eta\psi^{*}(\rho_{a}f)$, $\alpha_{a}(\psi^{*}(f)) = \eta\psi(\rho_{a}f)$. Since $\tilde J_{a}$ is antiunitary and fixes $\Omega$, $\omega(\alpha_{a}(A)) = \overline{\omega(A)}$. For $A = \psi^{*}(f)\psi(g)$: $\omega(\alpha_{a}(A)) = \omega(\psi(\rho_{a}f)\psi^{*}(\rho_{a}g)) = \ip{f}{\rho_{a}(1 - T)\rho_{a}g}$ and $\overline{\omega(A)} = \ip{f}{Tg}$, so
\[ \rho_{a}(1 - T)\rho_{a} = T. \tag{1} \]
Also $J_{a}\mathcal{R}_{a}J_{a} = Z\mathcal{R}_{a}^{c}Z^{*}$, since $J_{a}$ maps the generators $\pi(\psi(f))$, $f \in h_{a}$, of $\mathcal{R}_{a}$ onto $\eta Z\pi(\psi^{*}(\rho_{a}f))Z^{*}$, which generate $Z\mathcal{R}_{a}^{c}Z^{*}$.

\emph{Step 2 (purity, clause (i)).} By Tomita, $\mathcal{R}_{a}' = J_{a}\mathcal{R}_{a}J_{a} = Z\mathcal{R}_{a}^{c}Z^{*}$ (twisted Haag duality at the cut). $\Omega$ is cyclic for $\mathcal{R}_{a}$, so $\mathcal{R}_{a}$ on $\mathcal{H}$ is the GNS von Neumann algebra of the restriction of $\omega$ to $\CAR(h_{a})$, a gauge-invariant quasi-free state, hence a factor (Fact \claimref{synthesis.imported-powers-stormer-takesaki-disjointness}(a)); likewise $\mathcal{M} = \pi(\CAR(h))''$ is a factor. Let $x \in \mathcal{M}'$; $\mathcal{M}'$ is $\Theta$-invariant, so split $x = x_{+} + x_{-}$ into even and odd parts, both in $\mathcal{M}'$. The even part commutes with $\mathcal{R}_{a}^{c}$ and with $\Theta$, hence with every $ZYZ^{*} = Y_{+} + I_{4}\Theta Y_{-}$, $Y \in \mathcal{R}_{a}^{c}$: $x_{+} \in (Z\mathcal{R}_{a}^{c}Z^{*})' = \mathcal{R}_{a}$, and $x_{+} \in \mathcal{M}' \subset \mathcal{R}_{a}'$, so $x_{+}$ is central in the factor $\mathcal{R}_{a}$, a scalar. So the even part of $\mathcal{M}'$ is $\C 1$. If $v \ne 0$ is odd in $\mathcal{M}'$, then $vv^{*}$ and $v^{*}w$ (for odd $w \in \mathcal{M}'$) are even, hence scalars, $vv^{*} = c > 0$, and $w = c^{-1}vv^{*}w = c^{-1}v(v^{*}w) \in \C v$; with $v^{2} \in \C$ this makes $\mathcal{M}' = \mathrm{span}\{1, v\}$ abelian, so $\mathcal{M}' \subset \mathcal{M}'' = \mathcal{M}$ and $\mathcal{M}' \subset \mathcal{M}\cap\mathcal{M}' = \C 1$, contradicting oddness of $v \ne 0$. Hence $\mathcal{M}' = \C 1$, $\pi$ is irreducible and $\omega$ is pure, so $T$ is a projection. This last step is the only-if direction of Araki's purity criterion (a gauge-invariant quasi-free state is pure if and only if its symbol is a projection; Fact \claimref{synthesis.imported-powers-stormer-takesaki-disjointness}(d) \cite{Araki1970}); the factor property used above is clause (a) of the same Fact. States KMS for the affine translations at positive temperature, and the stationary half-thermal products other than $\omega_{0}$, are not pure ($T(1 - T) \ne 0$), so they satisfy CGMA$_{\eta}$ at no cut.

\emph{Step 3 (stationarity and charge symmetry, clause (ii)).} From (1) at $a$ and at $b$, $T = \rho_{a}(1 - T)\rho_{a} = \rho_{a}\rho_{b}T\rho_{b}\rho_{a}$, and $\rho_{a}\rho_{b}f(u) = f(u - 2(a - b))$ is the translation by $2(a - b)$ (checked as affine maps). So $T$ commutes with the translations by the subgroup generated by $2(A - A)$; if it is dense, strong continuity of the translation group on $h$ gives commutation with every translation. An operator on $L^{2}(\R)\otimes K$ commuting with all translations is a bounded measurable $B(K)$-valued function of $P$: $T = \theta(P)$. Since $\rho_{0}P\rho_{0} = -P$, (1) reads $\theta(-p) = 1 - \theta(p)$ for almost every $p$. Steps 1 and 3 use only the state-level identity (1); no positivity and no modular group.

\emph{Step 4 (the natural-cone form).} Fix the cut $0$ (translate otherwise; $T$ is translation invariant). For $x \in h_{0}$, $B(x)$ is self-adjoint in $\mathcal{R}_{0}$ and odd, and by Step 1, $J_{0}\pi(B(x))J_{0} = \eta Z\pi(B(\rho_{0}x))Z^{*}$, so $J_{0}\pi(B(x))\Omega = \eta Z\pi(B(\rho_{0}x))\Omega = -I_{4}\eta\,\pi(B(\rho_{0}x))\Omega$. Step 0 gives
\[ 0 \le \ip{B(x)\Omega}{J_{0}B(x)\Omega} = -I_{4}\eta\,\omega(B(x)B(\rho_{0}x)). \]
Now $\ip{x}{\rho_{0}x} = 0$ (disjoint supports) and $\rho_{0}T = (1 - T)\rho_{0}$ by (1), so $\ip{\rho_{0}x}{Tx} = \ip{x}{\rho_{0}Tx} = -\ip{x}{T\rho_{0}x}$ and $\omega(B(x)B(\rho_{0}x)) = -2\ip{x}{T\rho_{0}x}$. Take $x = \sum_{k}e^{-a_{k}u}\xi_{k}$ ($a_{k} > 0$, $\xi_{k} \in K$), dense in $h_{0}$ (a function orthogonal to every $e^{-au}$ has vanishing Laplace transform); $\hat x(p) = \sum_{k}\xi_{k}/(a_{k} - I_{4}p)$. Splitting $\int_{\R}$ into $p > 0$ and $p < 0$ and using $\theta(-p) = 1 - \theta(p)$, with $c_{kl}(p) = 1/((a_{k} - I_{4}p)(a_{l} - I_{4}p))$,
\[ \ip{x}{T\rho_{0}x} = \frac{1}{2\pi}\sum_{k, l}\int_{0}^{\infty}\Big(\ip{\xi_{k}}{\xi_{l}}c_{kl} - 2I_{4}\ip{\xi_{k}}{\theta(p)\xi_{l}}\,\mathrm{Im}\,c_{kl}\Big)dp. \]
Here $\mathrm{Im}\,c_{kl} = \kappa_{p}(a_{k}, a_{l})$ with $\kappa_{p}(a, b) = p(a + b)/((a^{2} + p^{2})(b^{2} + p^{2}))$, the double Laplace transform of $\sin(p(u + w))$, and $\int_{0}^{\infty}c_{kl}\,dp = I_{4}\ln(a_{l}/a_{k})/(a_{l} - a_{k}) = I_{4}\int_{0}^{\infty}\kappa_{p}(a_{k}, a_{l})\,dp$ (the real part integrates to $0$; for $a = b$, $\int_{0}^{\infty}(a^{2} - p^{2})/(a^{2} + p^{2})^{2}dp = [p/(a^{2} + p^{2})]_{0}^{\infty} = 0$). Hence $\ip{x}{T\rho_{0}x}$ is purely imaginary and the positivity reads
\[ \sum_{k, l}\int_{0}^{\infty}\ip{\xi_{k}}{w(p)\xi_{l}}\,\kappa_{p}(a_{k}, a_{l})\,dp \ \ge\ 0, \qquad w(p) = -\eta\,(1 - 2\theta(p))\ (p > 0). \tag{2} \]
For the null vacuum and $\eta = -1$, $w = 1$ and the form on $x_{m} = u^{m}e^{-u}$ is $\tfrac{1}{\pi}\cdot\tfrac{m!\,n!}{m + n + 1}$ (the kernel $\int_{0}^{\infty}\sin(ps)\,dp = 1/s$; Beta integral times $(m + n)!$), a positive matrix; with $\eta = +1$ it is negative: positivity fixes $\eta = -1$ for $\omega_{0}$ and $\eta = +1$ for its time reverse (checked exactly).

\emph{Step 5 (Loewner, clause (iii)).} Fix a unit $\xi \in K$, put $w_{\xi}(p) = \ip{\xi}{w(p)\xi} \in [-1, 1]$, $g_{p}(a) = p/(a^{2} + p^{2})$ and $H(a) = -\int_{0}^{\infty}w_{\xi}(p)\,(g_{p}(a) - g_{p}(1))\,dp$ (absolutely convergent, locally uniformly for $\mathrm{Re}\,z > 0$ with $a \to z$, since $g_{p}(z) - g_{p}(1) = O(p^{-3})$ at large $p$; the poles of $g_{p}(z)$ are at $z = \pm I_{4}p$). The identity $\kappa_{p}(a, b) = (g_{p}(a) - g_{p}(b))/(b - a)$ (checked exactly) turns (2) with all $\xi_{k} = \xi$ into positivity of the Loewner matrices $[(H(a_{k}) - H(a_{l}))/(a_{k} - a_{l})]$, diagonal $H'(a_{k})$, for every finite set of points of $(0, \infty)$. By Loewner's theorem (Fact \claimref{synthesis.imported-loewner-theorem-and-yafaev-hankel-sign-criterion}(a) \cite{Loewner1934}) $H$ is the restriction to $(0, \infty)$ of a Pick function $F$, analytic on the upper half-plane with $\mathrm{Im}\,F \ge 0$, which continues analytically across $(0, \infty)$ and by reflection to $\C\setminus(-\infty, 0]$. $H$ is analytic on $\{\mathrm{Re}\,z > 0\}$ and agrees with $F$ on $(0, \infty)$, so $F = H$ on the right half-plane. For $y > 0$ the point $I_{4}y$ lies in the upper half-plane, where $F$ is continuous, so $0 \le \mathrm{Im}\,F(I_{4}y) = \lim_{\delta\downarrow 0}\mathrm{Im}\,H(\delta + I_{4}y)$. With $z = \delta + I_{4}y$, $z^{2} + p^{2} = (p^{2} - y^{2} + \delta^{2}) + 2I_{4}\delta y$ and
\[ \mathrm{Im}\,H(\delta + I_{4}y) = \int_{0}^{\infty}w_{\xi}(p)\,P_{\delta}(p; y)\,dp, \qquad P_{\delta}(p; y) = \frac{2\delta y\,p}{(p^{2} - y^{2} + \delta^{2})^{2} + 4\delta^{2}y^{2}} \ge 0. \]
In $t = p^{2}$, $\int_{0}^{\infty}P_{\delta}\,dp = \tfrac12\big(\tfrac{\pi}{2} + \arctan\tfrac{y^{2} - \delta^{2}}{2\delta y}\big) \to \tfrac{\pi}{2}$, the mass outside any neighbourhood of $p = y$ tends to $0$, and near $p = y$, $P_{\delta}$ is bounded by a constant multiple of the Poisson kernel $\delta/((p - y)^{2} + \delta^{2})$; so at every Lebesgue point of $w_{\xi}$, $\mathrm{Im}\,F(I_{4}y) = \tfrac{\pi}{2}w_{\xi}(y)$ (the standard Lebesgue-point lemma for Poisson-type kernels; argued at that level). Hence $w_{\xi} \ge 0$ almost everywhere, and over a countable dense set of $\xi$, $w(p) \ge 0$ for almost every $p > 0$: $\theta(p) \le \tfrac12$ for $p > 0$ when $\eta = -1$, $\theta(p) \ge \tfrac12$ when $\eta = +1$. Conversely, for $K$ one-dimensional, (2) for all finite families is exactly the operator monotonicity of $H$ on $(0, \infty)$ (the exponentials being dense and the form bounded); it is strictly stronger than $w \ge 0$ (for $w = \chi_{(1, 2)}$, $H = -\tfrac12\ln\frac{a^{2} + 4}{a^{2} + 1} + \mathrm{const}$ has branch points at $I_{4}$ and $2I_{4}$ and is not a Pick function). The same conclusion follows from the Laplace representation of nonnegative Hankel operators (Fact \claimref{synthesis.imported-loewner-theorem-and-yafaev-hankel-sign-criterion}(b); \cite{Yafaev2015}, Theorem 4.7 of the arXiv version, Theorem 4.9 of the journal version): the Hankel kernel of (2) is $k(s) = \int_{0}^{\infty}w_{\xi}(p)\sin(ps)\,dp$, and $k = \int e^{-s\lambda}d\mu(\lambda)$ with $\mu \ge 0$ makes the Gram kernel $\int d\mu(\lambda)/((a + \lambda)(b + \lambda))$ and $H$ a Stieltjes function, visibly Pick.

\emph{Step 6 (the ground-state clause, clause (iv)).} By Step 2, $T = \theta(P)$ is a projection, so $w(p) = -\eta(1 - 2\theta(p))$ is a self-adjoint unitary for each $p > 0$; $w \ge 0$ forces $w = 1$, $1 - 2\theta(p) = -\eta$, so $\theta(p) = \tfrac12(1 + \eta)\id_{K}$ and $\theta(-p) = \tfrac12(1 - \eta)\id_{K}$ for almost every $p > 0$. For $\eta = -1$: $T = \chi_{(-\infty, 0)}(P)\otimes\id_{K}$, the null vacuum, the gauge-invariant quasi-free ground state of the affine translations, with positive generator (Theorem \claimref{synthesis.null-vacuum-on-a-horizon-generator}(ii)). For $\eta = +1$: $T = \chi_{(0, \infty)}(P)\otimes\id_{K} = \rho_{0}T_{0}\rho_{0}$, the time reverse. No stationarity was assumed (it came out in Step 3), no temperature and no modular group (only $J_{a}$, through $J_{a}\mathcal{R}_{a}J_{a} = \mathcal{R}_{a}'$ and Step 0).

\emph{Step 7 (the converse, imported from memory).} The pinned Fact \claimref{angle.imported-theorems-deliver-their-hypotheses} (its abstract form for half-line algebras with positive translations; Borchers' $JU(s)J = U(-s)$ is also in Fact \claimref{synthesis.imported-kuckert-modular-pct-spin-statistics-and-borchers-theorem}(ii)) and the stored theorem \claimref{synthesis.thermality-of-every-future-half-generator-fixes-the-period-and-the-null-ground-state}(ii), by which $\omega_{0}$ is KMS at $2\pi$ for the boost about every point, give the modular group: $\Delta_{a}^{I_{4}t}$ is the boost about $a$ at rapidity $-2\pi t$. That the modular conjugation acts on the fields as the twisted charge-conjugating reflection about $a$ is not in that Fact; it is Fact \claimref{synthesis.imported-modular-conjugation-of-the-chiral-fermi-field-is-the-twisted-reflection}(c), resting on \cite{BGL1993} and, for the twisted fermionic second quantisation, \cite{BJL2002}, the field-level form cited there from memory. On one-particle vectors this is the analytic continuation of the boost to rapidity $I_{4}\pi$: for $x, y \in h_{a}$, $\ip{B(y)\Omega}{\Delta_{a}^{1/2}B(x)\Omega}$ continues $\omega(B(y)\beta^{(a)}_{-2\pi t}(B(x)))$ to $t = -I_{4}/2$, where $\psi(u) \mapsto e^{\theta/2}\psi(a + e^{\theta}(u - a))$ at $\theta = I_{4}\pi$ gives $I_{4}\psi(2a - u)$, and the vacuum kernel $\tfrac{I_{4}}{2\pi}(u - v + I_{4}0)^{-1}$ is analytic along the path; so $\ip{B(y)\Omega}{J_{a}B(x)\Omega} = I_{4}\,\omega(B(y)B(\rho_{a}x))$, which is the value $-I_{4}\eta\,\omega(B(y)B(\rho_{a}x))$ of Step 4 with $\eta = -1$ (argued); the sign so obtained agrees with the one fixed independently by the proved positivity of Step 4, a consistency check, not a proof of the converse. So $\omega_{0}$ satisfies CGMA$_{-1}$ at every cut and its time reverse CGMA$_{+1}$; with Step 6, under the assumption $\eta = -1$ (the sign of the unit $I_{4}$ relative to time being a datum the series does not fix), CGMA$_{-1}$ at the cuts of $A$ is equivalent to (N$_{+}$) on the generator.

\emph{Step 8 (sharpness, clause (vi)).} (a) For a measurable unitary-valued $g$ on $\R$ (a local gauge transformation, $\gamma_{g}(\psi(f)) = \psi(gf)$), $\omega_{0}\circ\gamma_{g}$ has the same GNS space, vector and algebras $\mathcal{R}_{a}$ as $\omega_{0}$ ($\gamma_{g}$ preserves every $h_{I}$), hence the same $J_{a}$; in its representation $J_{a}$ acts as $\psi(f) \mapsto -Z\psi^{*}(\overline{g}\,(g\circ\rho_{a})\,\rho_{a}f)Z^{*}$, which is CGMA$_{-1}$ at $a$ exactly when $g\circ\rho_{a} = g$. With $g$ equal to $\tfrac{3 + 4I_{4}}{5}$ where the distance from $u$ to $2\Z$ is below $\tfrac12$ and to $1$ elsewhere, $g$ is symmetric about $0$ and $1$ (period $2$), so CGMA$_{-1}$ holds at every integer cut, while the symbol $\overline{g}T_{0}g$ has kernel $g(u)\overline{g(v)}\tfrac{I_{4}}{2\pi}(u - v)^{-1}$, invariant under translation by $2$ and not by $\tfrac12$ (exact values in the check): the state is not stationary, not $\omega_{0}$, so two cuts with commensurable distance do not suffice. (b) The translation-KMS states, $\theta(p) = (1 + e^{\beta p})^{-1}$, satisfy (1) at every cut, and their $w = \tanh(\beta p/2) \ge 0$; their Loewner function is $\psi(\beta a/2\pi + \tfrac12)$ up to a constant ($\psi$ here the digamma function, not the field) (Binet's integral, cited from memory; computed numerically to $10^{-9}$), and $\psi(z + \tfrac12) = -\gamma + \sum_{n \ge 0}\big(\tfrac{1}{n + 1} - \tfrac{1}{n + z + 1/2}\big)$ is a Pick function, so they pass (2) at every cut (argued); they fail only Step 2. (c) The state with $\theta = \chi_{S}$, $S = (-\infty, -2)\cup(0, 2)$, is pure, translation invariant and satisfies (1); on $e^{-\lambda u}$ the form (2) with $\eta = -1$ is $\int_{0}^{\infty}(1 - 2\theta)\kappa_{p}(\lambda, \lambda)\,dp = -1/\lambda + 2\lambda/(\lambda^{2} + 4)$ (antiderivative $-\lambda/(\lambda^{2} + p^{2})$), equal to $-\tfrac35$ at $\lambda = 1$ and $\tfrac{3}{20}$ at $\lambda = 4$, so the form is indefinite for either $\eta$ and CGMA$_{\eta}$ fails at every cut; its kernel differs from $\omega_{0}$'s by $\tfrac{1}{2\pi}\big(\int_{0}^{2} - \int_{-2}^{0}\big)e^{-I_{4}ps}dp = -\tfrac{I_{4}}{\pi}(1 - \cos 2s)/s$, an entire function, so it passes the generator microlocal condition at every point. It is locally normal to $\omega_{0}$ on every bounded interval $O$: $\chi_{O}(T_{S} - T_{0})\chi_{O} = \chi_{O}\chi_{(0, 2)}(P)\chi_{O} - \chi_{O}\chi_{(-2, 0)}(P)\chi_{O}$ is a difference of positive trace-class operators (each is $X^{*}X$ with $X = \chi_{(0, 2)}(P)\chi_{O}$, resp. $\chi_{(-2, 0)}(P)\chi_{O}$, Hilbert-Schmidt since its kernel is bounded and $O$ is bounded), so the compressed symbols $A$, $B$, and $1 - A$, $1 - B$, differ by trace-class operators; the Powers-Stormer inequality (Fact \claimref{synthesis.imported-powers-stormer-takesaki-disjointness}(c)) makes $A^{1/2} - B^{1/2}$ and $(1 - A)^{1/2} - (1 - B)^{1/2}$ Hilbert-Schmidt, and criterion (a) gives quasi-equivalence of the restrictions to the algebra of $O$. (d) The net-level form (only $J_{a}\mathcal{R}_{I}J_{a} = Z\mathcal{R}_{\rho_{a}I}Z^{*}$) is satisfied by every $\omega_{0}\circ\gamma_{g}$, by the net-level identity of the companion result \claimref{synthesis.modular-data-on-a-generator-fix-neither-the-translation-lift-nor-the-sign-of-i}(i) (not reproved here; it is the first sentence of (a)), including the translation-invariant Fermi-level states $g = e^{I_{4}ku}$, symbol $\chi_{(-\infty, k)}(P)$, which are not ground states for $k \ne 0$ and fail (1) ($\theta(-p) + \theta(p) = 2$ for $0 < p < k$); so the elementwise form is needed, in line with the Minkowski statement that net-level geometric modular action alone does not give the spectrum condition (Fact \claimref{synthesis.imported-buchholz-florig-summers-geometric-modular-action-and-modular-stability}(e) \cite{BFS1999}).
